\subsection{Rank-one augmentation modules in a finite category}

\begin{lemma}[Affine representatives for balanced degree-two classes]
\label{C:category-affine-balanced}
Let $C$ be a finite strongly connected category with rectangular-band
minimal ideal.  Use the row and column sets $R_i,C_i$ and the
augmentation modules $V,W$ of Lemma~\ref{C:category-band-Tor}.
Suppose $\dim V=1$.  Then every class in $H^2(BC;\mathbb Q)$ has
an affine cocycle representative of the following form.  There is a
unique object $a$ with $|R_a|=2$, while every other $R_i$ is a
singleton.  Choose functions $a_i:R_i\to\mathbb Q$ and
$b_i:C_i\to\mathbb Q$, both zero away from object $a$, and with
chosen distinguished values zero.  For each arrow $f:i\to j$ there
are scalars $\chi_f,a_f,b_f$ such that
\begin{align*}
 a_j\sigma_f&=\chi_f a_i+a_f,&
 b_i\tau_f&=\chi_f b_j+b_f.
\end{align*}
Here $\sigma_f:R_i\to R_j$ and $\tau_f:C_j\to C_i$ are the
row and column actions, and $\chi_f=0$ unless $i=j=a$.
For a composable pair written in multiplication order,
\[
 u(f,g)=b_fa_g
\]
represents the prescribed degree-two class.
\end{lemma}
\begin{proof}
The comparison of Lemma~\ref{C:category-band-Tor} identifies
degree-two cohomology with the dual of $W\otimes_A V$.
Object idempotents show that a balanced functional is supported on
$W_a\otimes V_a$.  Choose $a_a$ to take the values zero and one
on the two row points.  Any balanced functional is then
$\lambda(w,v)=b_a(w)a_a(v)$ for a linear functional $b_a$ on
the column augmentation space; extend $b_a$ to a function on
$C_a$ with value zero at its distinguished point.

On $V_a$, an endomorphism acts by a scalar $\chi_f$.
Balancing says that its column action has the same scalar on
$b_a$.  For a cross arrow leaving $a$, balancing says that
$b_a\tau_f$ is constant.  For a cross arrow entering $a$, its
row source is a singleton, so $a_a\sigma_f$ is constant.
These observations give exactly the displayed affine equations,
with $\chi_f=0$ on all other arrows.  The maps on the one-dimensional
module $V$ show that $\chi_{fg}=\chi_f\chi_g$ whenever $fg$ is
defined.  Composition of the affine maps gives
\begin{equation}\label{C:category-affine-composition}
 a_{fg}=a_f+\chi_f a_g,\qquad
 b_{fg}=b_g+\chi_g b_f.
\end{equation}
Thus $u$ is a normalized category two-cocycle, by substitution into
$\delta u(f,g,h)=u(g,h)-u(fg,h)+u(f,gh)-u(f,g)$.

We verify the asserted identification with the balanced functional,
rather than only the cocycle equation.  Let $o$ be the object of a
chosen minimal idempotent $e$.  Choose row arrows
$u_r:o\to i$ and column arrows $v_c:i\to o$ as in
Lemma~\ref{C:finite-category-Rees}.  Since the corner group is
trivial, $v_cu_r=e$ for every such pair.  For augmentation vectors
$w=\sum_c w_cv_c$ and $v=\sum_r v_ru_r$ at object $i$, set
\[
 z(w,v)=\sum_{c,r}w_cv_r[\,v_c\mid u_r\,].
\]
This is a two-cycle: in
$\partial[v_c\mid u_r]=[u_r]-[e]+[v_c]$, every summand vanishes
after one of the two zero coefficient sums is taken.  For
$f:i\to j$, the boundary of
$\sum_{c,r}w_cv_r[v_c\mid f\mid u_r]$ is
$z(w,fv)-z(wf,v)$, so these cycles respect the balancing relation.
They are the inverse of the two augmentation connecting maps in
Lemma~\ref{C:category-band-Tor}: lift their final coefficient at
$o$ to $e\in Re$; the first boundary is the bar one-cycle with
coefficients $v$ and column entries $v_c$; lifting its other
coefficient to $e\in eR$ gives boundary $w\otimes v$.
In particular they represent every degree-two homology class.
Finally the minimal-ideal actions are constant, and hence
$a_{u_r}=a_i(r)$ and $b_{v_c}=b_i(c)$.  Therefore
\[
 \langle u,z(w,v)\rangle
 =\sum_{c,r}w_cv_r b_i(c)a_i(r)=\lambda(w,v).
\]
This proves the claimed cohomology representative.
\end{proof}

\begin{proposition}[The category rank-one cup obstruction]
\label{C:category-rank-one-cups}
Under the hypotheses of Lemma~\ref{C:category-band-Tor}, if
$\dim V\leq1$ or $\dim W\leq1$, then every product of two
rational degree-two classes is zero.  Thus a nonzero such product
requires
\[
 \sum_i r_i\geq |\operatorname{Ob}C|+2,
 \qquad
 \sum_i c_i\geq |\operatorname{Ob}C|+2.
\]
No finite-homotopy hypothesis is needed.
\end{proposition}
\begin{proof}
Dimension zero gives $H^2=0$.  Taking the opposite category
exchanges $V,W$ and preserves the assertion, so assume $\dim V=1$
and use Lemma~\ref{C:category-affine-balanced}.  Let $e_i$ be one
at object $a$ and zero elsewhere.  The offsets $a_f,b_f$ are
supported, respectively, on arrows with target $a$ and source $a$.
The zero class needs no argument, so suppose that the balanced
functional is nonzero; both scalar value sets are then nonconstant.
Whenever $\chi_f\ne0$, the induced affine maps permute the finite
nonconstant value sets of $a_a$ and $b_a$.  Their slopes are
$1$ or $-1$.  A slope-one permutation has zero translation, and
all slope-minus-one permutations of one finite value set have
its same midpoint.  Choose rational midpoints $A,B$ for the two
sets, using zero if no negative-slope map occurs, and put
\[
 a'_f=a_f+A\chi_f-Ae_{t(f)},\qquad
 b'_f=b_f+B\chi_f-Be_{s(f)}.
\]
Then $\chi_fa'_f=\chi_fb'_f=0$ for every arrow.  The modified
offsets still satisfy \eqref{C:category-affine-composition},
because $\chi_f$ is supported on endomorphisms of $a$.

Set $u'(f,g)=b'_fa'_g$ and
$h(f)=(e_{s(f)}+e_{t(f)})/2$.  For a composable pair,
$\delta h(f,g)=e_{s(f)}=e_{t(g)}$.  The support of the original
offsets and $\chi_f$ gives the exact identity
\[
 u'-u=-\delta\bigl(A b+B a+AB(\chi-h)\bigr).
\]
The cochain in parentheses vanishes on every identity arrow.
Finally set
\[
 w(f,g,h)=b'_fa'_gb'_ga'_h.
\]
Expanding its coboundary with the affine composition equations
gives, on every composable quadruple,
\[
 \delta w(f,g,h,k)
 =u'(f,g)u'(h,k)
   +b'_fa'_k\,\chi_g b'_g\,\chi_h a'_h
 =u'(f,g)u'(h,k).
\]
Every degree-two square is therefore zero.  Over $\mathbb Q$,
graded commutativity and the square of a sum show that every
product of two degree-two classes is zero as well.
\end{proof}

\begin{lemma}[Rational Tor comparison with a finite minimal group]
\label{C:category-group-Tor}
The Tor formula of Lemma~\ref{C:category-band-Tor} remains valid
when the minimal group $G$ is any finite group, if $V,W$ are
defined by the orbit augmentations
\[
 P=Re\otimes_{\mathbb QG}\mathbb Q\longrightarrow\mathbf1_\ell,
 \qquad
 Q=\mathbb Q\otimes_{\mathbb QG}eR\longrightarrow\mathbf1_r.
\]
Their dimensions are still $\sum_i(r_i-1)$ and
$\sum_i(c_i-1)$, and $A=R/\mathbb QK$.
\end{lemma}
\begin{proof}
The trivial $\mathbb QG$-module is projective, by the averaging
idempotent $|G|^{-1}\sum_{g\in G}g$.  Thus $P,Q$ are projective
left and right $R$-modules, respectively.  Their bases are the
row and column orbits, and their augmentation kernels are
annihilated by $K$.  The Rees multiplication bijection in
Lemma~\ref{C:finite-category-Rees} gives
\[
 Re\otimes_{\mathbb QG}eR\cong\mathbb QK.
\]
The row and column $G$-actions are free, so this ideal is
projective on each side.  Also $eV=We=0$ and
$Q\otimes_RV=\mathbb Q\otimes_{\mathbb QG}eV=0$.
The two augmentation sequences therefore give
$H_q(BC;\mathbb Q)\cong\operatorname{Tor}_{q-2}^{R}(W,V)$,
including $q=2$.  Finally
$W\otimes_R\mathbb QK=We\otimes_{\mathbb QG}eR=0$;
the projective resolution $0\to\mathbb QK\to R\to A\to0$
and the same change-of-rings argument as before replace $R$
by $A$.
\end{proof}

\begin{theorem}[The strongly connected two-object exclusion through twenty]
\label{C:twenty-strongly-connected}
A finite category with two objects, nonempty hom-sets in both
directions, and at most twenty arrows cannot have classifying
space homotopy equivalent to $\mathbb{CP}^2$.
\end{theorem}
\begin{proof}
If either object identity is in the minimal ideal, the retract
deletion in Lemma~\ref{C:finite-category-Rees} leaves an
equivalent one-object category with at most nineteen arrows.
This contradicts Corollary~\ref{B:CP2-lower21}.  Thus both
identities are outside $K$.

If all row numbers or all column numbers are one, the same
Rees lemma gives $BC\simeq BG$, which has no positive rational
homology.  Otherwise both total row and column numbers are at
least three.  If $G\ne1$, the arrow bound
\[
 |G|\Bigl(\sum_i r_i\Bigr)\Bigl(\sum_i c_i\Bigr)+2\leq20
\]
forces $|G|=2$, both sums equal to three, and exactly two arrows
outside $K$, namely the identities.  Then $A\cong\mathbb Q^2$
is semisimple, so Lemma~\ref{C:category-group-Tor} gives
$H_4(BC;\mathbb Q)=0$.

It remains to treat $G=1$.  The nonzero degree-two square of
$\mathbb{CP}^2$ and Proposition~\ref{C:category-rank-one-cups}
force both row and column sums to be at least four.  Thus
$|K|\geq16$ and $\dim A\leq4$.  The two object identities
survive as nonzero orthogonal idempotents in $A$, so $A$ is
nonlocal.  By Lemma~\ref{B:nonlocal-four-Tor}, a nonzero
$\operatorname{Tor}_2^A(W,V)$ forces nonzero Tor in arbitrarily
high degrees.  Lemma~\ref{C:category-band-Tor} would therefore
make nonzero $H_4$ incompatible with eventual vanishing of
$H_q(BC;\mathbb Q)$.  This is the required contradiction.
\end{proof}
